\section{Experiments}
\label{sec:experiments}

\subsection{Dataset}
\label{sec:dataset}

We use the MAVIC-T 2025 dataset from the challenge organizers, with imagery from UNICORN, USGS EROS HRO, and UMBRA across four regions: Bingham, Centerfield, Manhattan, and UC Davis. SAR$\rightarrow$EO operates at $256{\times}256$; the remaining three tasks use $1024{\times}1024$. For SAR$\rightarrow$IR and SAR$\rightarrow$RGB, we spatially align SAR scenes with optical mosaics before resampling. We create crop-augmented training sets for the $1024{\times}1024$ tasks (sliding windows, stride 512) to increase data diversity.

\begin{table}[h]
\centering
\caption{Task specifications and training data.}
\label{tab:dataset}
\resizebox{\columnwidth}{!}{%
\begin{tabular}{lcccc}
\toprule
Task & Res. & Channels & Orig. & Aug. \\
\midrule
SAR$\rightarrow$EO  & $256^2$  & $1{\to}1$ & 89,411 & -- \\
SAR$\rightarrow$RGB & $1024^2$ & $1{\to}3$ & 10,576 & 75,703 \\
SAR$\rightarrow$IR  & $1024^2$ & $1{\to}1$ & 10,576 & 74,317 \\
RGB$\rightarrow$IR  & $1024^2$ & $3{\to}1$ &  2,272 &  7,388 \\
\bottomrule
\end{tabular}
}%
\end{table}

\subsection{Training Details}
\label{sec:training}

All DBIM models are trained with the Prodigy adaptive optimizer~\cite{mishchenkoProdigyExpeditiouslyAdaptive2024}. For $1024{\times}1024$ tasks, we train at $512{\times}512$ and infer at full resolution (resolution curriculum). The bridge follows a VP log-SNR schedule with $\beta_d=2.0$, $\beta_{\min}=0.1$.

\begin{table}[h]
\centering
\caption{Per-task DBIM configuration.}
\label{tab:arch}
\resizebox{\columnwidth}{!}{%
\begin{tabular}{lcccc}
\toprule
Parameter & S2EO & S2RGB & S2IR & R2IR \\
\midrule
Base channels & 64 & 96 & 96 & 128 \\
Channel mults & (1,2,3,4) & \multicolumn{3}{c}{(1,1,2,2,4,4)} \\
Attention res. & N/A & \multicolumn{3}{c}{32, 16, 8} \\
Training res. & $256^2$ & \multicolumn{3}{c}{$512^2$} \\
Epochs & 55 & 24 & 39 & 86 \\
\bottomrule
\end{tabular}
}%
\end{table}

\subsection{Quantitative Results}
\label{sec:quantitative}

\cref{tab:main_results} reports the performance of our \textbf{EarthBridge} final submission (DBIM for SAR$\rightarrow$EO, SAR$\rightarrow$RGB, RGB$\rightarrow$IR; CUT for SAR$\rightarrow$IR) and selected baselines. The composite task score is defined as:
\begin{equation}
    S_{\text{task}} = \frac{\tfrac{2}{\pi}\arctan(\text{FID}) + \text{LPIPS} + \text{L1}}{3},
\end{equation}
where the $\arctan$ normalization maps FID to $[0,1]$.

\begin{table}[h]
\centering
\caption{Results on the MAVIC-T evaluation set. $\downarrow$ lower is better. \textbf{Bold} marks our final submission per track.}
\label{tab:main_results}
\resizebox{\columnwidth}{!}{%
\begin{tabular}{llcccc}
\toprule
Track & Method & FID$\downarrow$ & LPIPS$\downarrow$ & L1$\downarrow$ & Score$\downarrow$ \\
\midrule
\multirow{2}{*}{SAR$\rightarrow$EO}
  & DBIM (100 steps)  & 0.451 & 0.477 & 0.077 & 0.335 \\
  & \textbf{DBIM (500 steps)} & \textbf{0.222} & \textbf{0.504} & \textbf{0.080} & \textbf{0.269} \\
\midrule
\multirow{2}{*}{SAR$\rightarrow$RGB}
  & DDBM (100 steps) & 0.894 & 0.656 & 0.236 & 0.595 \\
  & \textbf{DBIM (1000 steps)} & \textbf{0.878} & \textbf{0.636} & \textbf{0.214} & \textbf{0.576} \\
\midrule
\multirow{2}{*}{SAR$\rightarrow$IR}
  & DBIM (1000 steps) & 0.958 & 0.630 & 0.194 & 0.594 \\
  & \textbf{CUT (huge)} & \textbf{0.645} & \textbf{0.603} & \textbf{0.145} & \textbf{0.464} \\
\midrule
\multirow{3}{*}{RGB$\rightarrow$IR}
  & DBIM (1 step)    & 0.786 & 0.293 & 0.093 & 0.390 \\
  & DBIM (10 steps)  & 0.380 & 0.150 & 0.090 & 0.206 \\
  & \textbf{DBIM (5 steps)}  & \textbf{0.357} & \textbf{0.151} & \textbf{0.090} & \textbf{0.200} \\
\bottomrule
\end{tabular}
}%
\end{table}

\paragraph{Key Observations.}
SAR$\rightarrow$EO and RGB$\rightarrow$IR are easier tasks: the former benefits from shared imaging geometry, the latter from RGB's rich three-channel input. SAR$\rightarrow$RGB and SAR$\rightarrow$IR are more challenging due to the large modality gap. For RGB$\rightarrow$IR, even one inference step yields reasonable results, but 5 steps achieves the best composite task score (0.200 vs.\ 0.206 at 10 steps), balancing FID, LPIPS, and L1. For SAR$\rightarrow$IR, CUT outperforms DBIM by a large margin (score 0.464 vs.\ 0.594), demonstrating that contrastive unpaired methods can be more effective when modality gaps are severe.

\subsection{Qualitative Results}
\label{sec:qualitative}

\begin{figure}[h]
    \centering
    \fbox{\rule{0pt}{2in} \rule{0.95\columnwidth}{0pt}}
    \caption{Qualitative results placeholder. Rows: SAR$\rightarrow$EO, SAR$\rightarrow$RGB, SAR$\rightarrow$IR, RGB$\rightarrow$IR. Columns: source image, EarthBridge output, ground-truth target.}
    \label{fig:qualitative}
\end{figure}

\cref{fig:qualitative} shows representative outputs from all four tasks. EarthBridge preserves fine-grained structural information from the source modality while accurately synthesizing target-modality textures across both $256{\times}256$ and $1024{\times}1024$ resolutions.

\subsection{Inference Speed}
\label{sec:inference_speed}

We report single-image inference times on a single NVIDIA 4090 (24~GB): RGB$\rightarrow$IR DBIM (5 steps)~$\approx$1.5\,s; SAR$\rightarrow$EO DBIM (500 steps)~$\approx$0.42\,s; SAR$\rightarrow$IR CUT~$\approx$0.47\,s; SAR$\rightarrow$RGB DBIM (1000 steps)~$\approx$160\,s.
